%--------------------------------------------------------------------------------
%
%--------------------------------------------------------------------------------
\unnumberedChapter{Summary}

This part of the book introduced the hardware foundation used throughout the 
remainder of the LCS hardware design. The focus was not on building one specific
application module, but rather on creating a reusable controller platform
that can serve as the basis for many different LCS nodes.
The main controller board combines all the common elements required by an LCS
node:

\begin{smallGapItemize}
\item a powerful controller platform,
\item the CAN bus communication interface,
\item non-volatile storage,
\item power management,
\item standardized extension interfaces.
\end{smallGapItemize}

The board is therefore much more than a controller carrier. It is a complete
LCS node capable of processing messages, running the LCS runtime library and
interacting with other nodes on the LCS bus.

During the development of the LCS runtime library, this board served as the
primary development platform. It allowed the individual runtime functions,
message handling and configuration concepts to be developed and tested on
real hardware. However, its purpose goes far beyond being a development
vehicle. The same board can be used to build useful LCS modules by adding
dedicated extension hardware.

The extension connector is the key element that enables this flexibility.
Extension boards can add application specific functionality while the
controller foundation remains unchanged. A turnout controller, signal module,
sensor interface or other specialized node can reuse the same controller
platform and concentrate only on the hardware and firmware required for its
specific purpose.

The hardware schematics introduced throughout this chapter follow the same
building block philosophy. Communication interfaces, power circuits,
controllers, level shifters and other functional blocks are designed as
reusable elements. These building blocks will appear again in the following
hardware designs and form the foundation for more complex modules.

For specialized nodes, such as the base station or the block controller, the
same principles still apply. In these cases, the main controller and the
required extension functionality may be combined into one integrated PCB.
This approach reduces complexity for more demanding hardware modules while
keeping the same design concepts. Even these integrated designs retain the
external I2C extension capability where socket boards are useful for local
functions.

Before moving on to these specialized modules, one more aspect is worth
mentioning. Although all LCS nodes use the same runtime library concepts,
the underlying hardware resources are not necessarily identical. Different
controller families, board layouts or integrated hardware designs may assign
controller pins and peripherals differently.

A common approach in many embedded projects is to handle these differences
with a growing number of conditional compilation statements. Over time,
these \textbf{ifdef} sections often spread throughout the entire firmware,
making the software increasingly difficult to maintain and frankly 
to understand.

The \textbf{Controller Dependent Code (CDC)} layer introduced earlier avoids
this problem. Hardware-specific details are isolated in a dedicated layer
that maps symbolic hardware resources used by the runtime to the actual
controller implementation. This keeps the LCS runtime portable while still
allowing optimized hardware designs.

With the concepts, firmware architecture and hardware foundation now in
place, we are ready to look at the first specialized LCS module. The next
chapter introduces the heart of the layout control system: the LCS base
station.
